% Generated draft: requires mathematical and source review.
\documentclass[11pt]{article}
\usepackage[margin=25mm]{geometry}
\usepackage{fontspec}
\setmainfont{Latin Modern Roman}
\usepackage{amsmath,amssymb,mathtools,unicode-math}
\setmathfont{Latin Modern Math}
\usepackage{enumitem,needspace,xurl,hyperref}
\setlength{\parindent}{0pt}
\setlength{\parskip}{5pt}
\setlength{\emergencystretch}{3em}
\newcommand{\N}{\mathbb{N}}
\newcommand{\Z}{\mathbb{Z}}
\newcommand{\Q}{\mathbb{Q}}
\newcommand{\R}{\mathbb{R}}
\newcommand{\C}{\mathbb{C}}
\newcommand{\F}{\mathbb{F}}
\newcommand{\A}{\mathbb{A}}
\newcommand{\PP}{\mathbb{P}}
\newcommand{\E}{\mathbb{E}}
\newcommand{\eps}{\varepsilon}
\newcommand{\dd}{\,\mathrm{d}}
\begin{document}
\section*{UnsolvedMath ID: 4200014}
\textbf{Source problem code:} catalogue entry
\subsection{Definitions and mathematical statement}
Let \(\mathbb T^3=\mathbb R^3/\mathbb Z^3\), let \(g_0\) be a flat metric, and consider metrics \(g\) close to \(g_0\). The source asks whether quasiperiodic globally minimizing geodesics exist in this setting, but does not specify the topology of closeness or the precise quasiperiodicity requirement.

A critical distinction is the space in which minimization is tested. For the lifted Riemannian metric \(\widetilde g\) on \(\mathbb R^3\), a unit-speed complete geodesic \(\gamma:\mathbb R\to\mathbb R^3\) is globally minimizing if
\[
 \operatorname{length}_{\widetilde g}(\gamma|_{[s,t]})
 =d_{\widetilde g}(\gamma(s),\gamma(t))=t-s
 \qquad(s<t).
\]
If instead the source's ``segment between any two of its points'' were interpreted using unrestricted distance on the compact torus itself, arbitrarily long unit-speed segments could not all minimize because that distance is bounded. Thus minimization on the universal cover, or an equivalent specified homotopy-constrained notion downstairs, needs confirmation rather than silent insertion.

One standard phase-space meaning of quasiperiodicity is that the projected geodesic-flow orbit has the form \(F(\theta+t\omega)\), where \(F:\mathbb T^m\to S_g\mathbb T^3\) parametrizes an invariant torus and the components of \(\omega\in\mathbb R^m\) are rationally independent. Here \(S_g\mathbb T^3\) is the unit tangent bundle. The source does not specify \(m\), the regularity or injectivity of \(F\), whether periodic orbits count, or whether full three-frequency motion is required.

Under a common intended reading the question would ask whether there is a neighborhood \(\mathcal U\) of \(g_0\) such that every \(g\in\mathcal U\) has at least one complete geodesic whose lift is globally minimizing and whose projected flow is quasiperiodic in the prescribed sense. This is an explicitly conditional reconstruction: the excerpt does not fix the metric regularity, Riemannian versus more general metrics, neighborhood topology, or whether the assertion is for every nearby metric, for generic metrics, or for prescribed rotation data.

The supplied summary attributes classical existence of global minimizers and minimal measures on tori to Mather--Hedlund--Bangert, while saying that it does not cleanly resolve the specific quasiperiodic near-flat three-torus question. No exact supporting theorem references are supplied. The 2000 source's open label has no recorded later authoritative resolution check.
\subsection{Short English statement}
Must metrics sufficiently close to a flat three-torus have a quasiperiodic geodesic that is globally minimizing in the intended lifted sense? The source leaves the minimizing space, closeness topology, frequency requirement, and quantifiers unspecified; classical minimal-geodesic existence alone does not identify the requested target.
\subsection{Sources}
\begin{itemize}
\item Local registry: \url{data/unsolvedmath/problems.json}, permanent UnsolvedMath ID \(4200014\), source code catalogue entry; supplied statement, background, and research summary, updated 2026-08-06.
\item Knill, \emph{Open problems in Hamiltonian dynamics} (2000), item 14, ``Mather theory near integrable systems'': \url{https://people.math.harvard.edu/~knill/seminars/intr/}. HTML extraction accessed 2026-07-29; no later authoritative resolution check is recorded.
\item The research summary names Mather, Hedlund, and Bangert for classical minimizer/minimal-measure theory but supplies no original publications or theorem statements resolving this question.
\end{itemize}
\end{document}
